% capitolo 6: data quality e observability
\chapter{Data Quality e Data Observability}
\label{cap:quality-observability}

\intro{Il seguente capitolo illustra le funzionalità di Data Quality e gli strumenti di monitoraggio dei dati nativi di OpenMetadata.}\\

\section{Agenti automatici}
All'interno di OpenMetadata, la raccolta continuativa dei metadati e il monitoraggio dello stato di salute delle informazioni sono affidati al motore di automazione \textbf{AutoPilot}. 
Si tratta di un sistema di agenti automatici che orchestrano le operazioni di estrazione, profilazione, classificazione e tracciamento senza richiedere configurazioni 
manuali complesse o interventi ripetitivi da parte dell'utente.

Gli agenti si attivano automaticamente nel momento in cui viene registrato un nuovo servizio o quando viene avviata una sincronizzazione con le fonti originali, 
operando in background secondo una pianificazione personalizzabile. Ciascun agente assolve a compiti specifici e mirati all'interno della piattaforma:
\begin{itemize}
    \item \textbf{Metadata Agent}: Costituisce l'agente primario ed è prerequisito per il funzionamento dell'intera piattaforma. Si occupa dell'estrazione e della sincronizzazione dei metadati dalle sorgenti collegate.
    \item \textbf{Profiler Agent}: Raccoglie metriche e statistiche descrittive sui dati a livello di tabella e di singola colonna (per esempio, calcolando la percentuale di valori nulli, il numero di valori distinti, gli intervalli minimi e massimi). I risultati alimentano direttamente la scheda di \emph{Data Quality} di ciascun dataset e forniscono la base quantitativa su cui vengono valutati i test di qualità.
    \item \textbf{AutoClassification Agent}: Analizza i campioni di testo prelevati dal Profiler Agent per alimentare il motore di \emph{Entity Recognition} basato su tecniche NLP, identificando automaticamente le informazioni personali e sensibili e applicando i relativi tag di sicurezza.
    \item \textbf{Lineage Agent}: Analizza i log delle query SQL e le definizioni delle viste per ricostruire in modo automatico il grafo delle dipendenze tra i dati.
    \item \textbf{Usage Agent e agenti specializzati}: Raccolgono le statistiche di utilizzo delle tabelle per alimentare il punteggio di popolarità (\emph{Popularity Scoring}) e le \emph{Query Analytics}.
\end{itemize}

Ogni agente è responsabile di una fase specifica del ciclo di vita dei metadati e opera in sequenza logica: poiché l'esecuzione di alcuni agenti è prerequisito per l'attivazione di altri, 
è fondamentale che vengano eseguiti nell'ordine corretto.

Lo stato di ciascun agente è visibile all'interno della scheda \emph{Agents} nella pagina del servizio di riferimento. Da questa interfaccia, gli amministratori e gli utenti 
autorizzati possono controllare lo stato di avanzamento delle attività (\emph{Active}, \emph{Pending}, \emph{Success}, \emph{Failed}), ispezionare i log di esecuzione per il 
\emph{debugging} o forzare un'esecuzione manuale a seguito di modifiche strutturali apportate a uno specifico dataset o servizio sorgente.


\section{Test di qualità}

\section{Gestione degli incidenti}

\section{Sistemi di alerting e notifica}
% e anche data collaboration